\section{Architecture}

\subsection{Components}

\begin{center}
\small
\begin{tabular}{@{}p{0.24\linewidth}p{0.66\linewidth}@{}}
\toprule
Component & Responsibility \\
\midrule
\code{cmd/fsb} & Command line, startup, the launch URL, opening the browser, shutdown; the remembered port, background mode and finding a running \texttt{fsb} (\code{config.go}, \code{background.go}, \code{process.go}). \\
\code{internal/rules} & The rule dialect: parsing, matching, name folding, the built-in rules. Knows nothing of files. \\
\code{internal/guard} & The chokepoint: opening, listing, searching, reading heads, metadata, sniffed endpoints, archive listing, the hard-link index. The only code that touches the file system on behalf of a request. \\
\code{internal/httpguard} & The middleware that makes a local web server safe: host and origin checks, the session, the prefix, security headers. Knows nothing of files. \\
\code{internal/quicklook} & Runs Apple's \code{qlmanage} to draw a Quick Look picture of a private copy that the guard made. The only code that runs another program on a file. Knows nothing of rules. \\
\code{internal/server} & The endpoints. Turns guard results into HTTP, and guard errors into status codes. \\
\code{web} & The embedded front end, and the fixed page that renders Markdown. \\
\code{e2e} & Browser tests, and their drivers for Chrome and Safari. \\
\code{specs} & These documents. \\
\bottomrule
\end{tabular}
\end{center}

\subsection{Structure}

\begin{center}
\begin{tikzpicture}[
  box/.style={draw, rounded corners=2pt, minimum width=2.6cm, minimum height=0.8cm, align=center, font=\small},
  arrow/.style={-{Stealth[length=2mm]}, thick},
  node distance=0.9cm and 1.4cm]
  \node[box] (browser) {Browser\\interface};
  \node[box, right=of browser] (hg) {httpguard\\middleware};
  \node[box, right=of hg] (srv) {server\\endpoints};
  \node[box, right=of srv] (guard) {guard};
  \node[box, below=of guard] (rules) {rules};
  \node[box, below=of srv] (fs) {file system\\(the OS)};
  \draw[arrow] (browser) -- (hg);
  \draw[arrow] (hg) -- (srv);
  \draw[arrow] (srv) -- (guard);
  \draw[arrow] (guard) -- (rules);
  \draw[arrow] (guard) -- (fs);
\end{tikzpicture}
\end{center}

Dependencies point one way: \code{server} uses \code{guard} and \code{httpguard};
\code{server} also uses \code{quicklook}, which draws only the copy that
\code{guard.CopyForQuickLook} made; \code{guard} uses \code{rules};
\code{rules}, \code{httpguard} and \code{quicklook} use nothing of the program's
own.  \code{cmd/fsb} assembles them.

\subsection{Life of a request}

\begin{enumerate}
  \item The HTTP server accepts the connection on the loopback address.
  \item The access log records the method, path and status, never the query
        (which is where the launch token travels).
  \item The middleware sets the security headers, then checks, in order: the
        \code{Host} header; the method; the \code{Origin}, \code{Referer} and
        \code{Sec-Fetch-Site} headers; that the path is under the random prefix
        (which it strips); the launch-token exchange; and the session cookie.  A
        failure of any check is the same 403 (\fq{SEC-2} to \fq{SEC-8}).
  \item The endpoint reads its parameters and calls one guard operation.
  \item The guard canonicalizes the path, checks it against the deny rules before and
        after opening, opens the file, verifies that what was opened is what was
        checked, and decides the file's kind from its bytes where needed.
  \item The endpoint turns the result into a response, and every guard error into a
        status code in one function (\code{fail}), so that denied, missing and
        outside-the-roots paths cannot be told apart by accident.
\end{enumerate}

\subsection{Startup}

\code{cmd/fsb} resolves the home folder to its real path.  \code{--init},
\code{--stop} and \code{--status} act and exit here, and \code{--background} starts
its child here (which runs these same steps from the start).  It then loads the two
rule files (or their built-in defaults); \code{--show-deny} and \code{--show-ignore}
print them and exit here, so they show exactly what would be enforced.  It then builds the guard from the roots (served once each, compared by
identity) and rules, names the credential locations the guard must also protect,
starts the hard-link index building in the background, binds the loopback port
(\code{choosePort}, Section~\ref{sec:port}), prints the launch URL, and only then
serves.  A failure at any step before serving stops the program with a message.

\subsection{The port, and background mode}
\label{sec:port}

Browsers keep local storage per origin, and the origin includes the port, so a port
chosen afresh at each start would reset every preference (\fq{CFG-4}).
\code{choosePort} therefore reuses the port in \code{\~{}/.config/fsb/lastport}, and
writes it after a start that did not use \code{--port} (best effort, like the front
end's own writes to local storage).  The random port was never a security control:
the boundary is the 256-bit launch token, the \code{Host} and \code{Origin} checks, and
the random path prefix, which still changes at every start.  \code{--port} is
deliberately not remembered, so that a one-off choice made to avoid a conflict heals
back to the usual port, and its preferences, once the conflict is gone.  If the
remembered port is taken, \code{fsb} fails rather than silently moving (which would
look as if it had lost the preferences); \code{portInUseMessage} asks \code{lsof} and
\code{ps} which program holds it so that the message can name it.

\code{--background} re-executes the program with \code{FSB\_BACKGROUND\_CHILD=1}
in a new session (\code{setsid}), so the child has no controlling terminal.  The
child's standard output and error are a writer (\code{handoff}) that goes to a pipe
on file descriptor~3 until the child is serving, then to
\code{\~{}/.config/fsb/background.log}.  The parent copies the pipe to its own output
until a fixed ready marker arrives (success: it releases the child and exits) or the
pipe closes without one (the child failed and has printed why; nothing is left
running).  The child records its process number in \code{\~{}/.config/fsb/pid} and
removes it on shutdown if it still names itself.

\code{findRunning} locates a running \code{fsb} for \code{--stop}, \code{--status} and
the refusal of a second \code{--background}: first the process in the PID file, then
the process listening on the remembered port.  Either is accepted only if \code{ps}
belongs to the same user and its executable file (from \code{/proc} or \code{lsof}, not
the name \code{ps} shows, which on macOS is the process's own \code{argv[0]}) is named
\code{fsb} or as the running binary is, and a process from the PID file must also have
the start time recorded there.  So a stale PID file whose number was reused, or another
program on the port that calls itself \code{fsb}, is never signaled.  \texttt{fsb}'s own
files are opened without following symbolic links and set to mode \code{0600}.
